\documentclass[10pt,a4paper]{article}
\usepackage[margin=1.7cm]{geometry}
\usepackage{graphicx,booktabs,tabularx,amsmath,xcolor,hyperref,caption,subcaption,enumitem,placeins}
\captionsetup{font=small,labelfont=bf,skip=2pt}
\setlist{nosep,leftmargin=*}
\setlength{\parskip}{2pt}
\setlength{\textfloatsep}{8pt}
\hypersetup{colorlinks,linkcolor=black,citecolor=black,urlcolor=blue!50!black}
\newcommand{\fig}[1]{../results/figures/#1}
\newcommand{\fillin}[1]{\colorbox{yellow!40}{\textbf{[#1]}}}
\newcommand{\ok}{\textcolor{green!45!black}{\textbf{reproduced}}}
\newcommand{\partly}{\textcolor{orange!80!black}{\textbf{partly}}}
\newcommand{\no}{\textcolor{red!70!black}{\textbf{not reproduced}}}
\title{\vspace{-1.3cm}\textbf{Network Dissection of AlexNet and VGG16: an honest reproduction in PyTorch}}
\author{Interpretable Machine Learning --- Group Project \\ \fillin{Group members (names, roll numbers)}}
\date{\fillin{Submission date} \quad Code: \url{https://github.com/Carbonyx-AI/netdissect-alexnet-vgg}}
\begin{document}
\maketitle
\vspace{-0.7cm}
\begin{abstract}\noindent
We re-implemented Network Dissection \cite{bau2017} from scratch in PyTorch and repeated the paper's experiments for
AlexNet and VGG16, using the paper's own Caffe models (AlexNet-Places205, -Hybrid, -Places365, BVLC AlexNet,
VGG16-Places365/Hybrid1365, loaded with our own converter) next to torchvision models. We report five results.
\textbf{(1)} Interpretability is tied to the basis: a random rotation of conv5 of the paper's AlexNet-Places205 removes 84\% of
the unique detectors (paper: 80\%) while the output is unchanged. \textbf{(2)} Detectors move from colour/texture to objects
with depth. \textbf{(3)} VGG16 $>$ AlexNet, trained $\gg$ random, and Places $>$ ImageNet for the original models.
\textbf{(4)} In small-scale training on Imagenette, seeds converge and detectors emerge during training, but batch normalisation
does not lower interpretability, removing dropout has no effect, and widening conv5 helps only a little.
\textbf{(5)} The method passes a randomisation sanity check, but single-unit labels have almost no causal importance, and
whether groups of same-labelled units matter depends on the model.
\end{abstract}

\section{Introduction}
Network Dissection is a \emph{post-hoc, model-specific, global} method (Molnar \cite{molnar}, \emph{Learned Features}): every unit of a
CNN is scored as a segmenter of every concept of a densely labelled dataset, and the number of distinct concepts that some
unit detects measures how interpretable a layer is. The paper asks whether interpretable units are an accident of the
coordinate system (H1) or a special, axis-aligned basis found by training (H2), and how interpretability depends on
layer, architecture, training data and training regime. We implemented the method and repeated the paper's experiments on
these questions (Sec.~3.2--3.5, 3.7) for AlexNet and VGG16, the two architectures of our topic, and then asked whether the
method itself can be trusted. Expectations were written down before running the experiments (\texttt{report/EXPECTATIONS.md}).

\section{Datasets and models}
\textbf{Broden} (\texttt{broden1\_224}, 63{,}305 images, labels at 112\,px, 1{,}197 concepts: 584 objects, 234 parts,
468 scenes, 32 materials, 47 textures, 11 colours, exactly the paper's Table~1) is the dataset the method is defined on;
dissection needs pixel-level concept labels. The official server was offline, so we used the Internet Archive copy and
verified it byte-for-byte (SHA-1). \textbf{Imagenette} (10 ImageNet classes) is used for accuracy checks, the fidelity test
and the small training runs. \fillin{note instructor approval of Imagenette}

\textbf{Models.} The paper used Caffe models, so we wrote a reader for \texttt{.caffemodel}/\texttt{.prototxt} files
(own protobuf decoder; grouped convolutions, LRN, ceil pooling, BGR mean subtraction). Every converted model was checked
before use: BVLC AlexNet reaches 76.7\% 1000-way top-1 on Imagenette (torchvision AlexNet 72.0\% under the same
protocol); on Broden scene images AlexNet-Places205 reaches 62.0\% top-1 (205 classes), AlexNet-Places365 46.6\%,
VGG16-Places365 51.5\%. AlexNet-Hybrid is weaker than expected on scenes (38.9\% on its 205 scene outputs) and is used with
that caveat. Not available: VGG16-Places205 (link dead) and the self-supervised AlexNets (server offline).

\textbf{Libraries.} PyTorch (computation, autograd), torchvision (AlexNet/VGG16 definitions and ImageNet weights),
NumPy/pandas (data handling), Pillow (images, label PNGs), SciPy (Schur decomposition, statistical tests), Matplotlib
(figures). Dissection, the Caffe reader, rotations and permutation importance are our code.

\section{Method}
For each unit $k$: $T_k$ is the exact top quantile with $P(a_k>T_k)=0.005$ over all positions of all images;
$A_k(x)$ is upsampled bilinearly to $112\times112$ with each activation at its receptive-field centre and thresholded,
$M_k=S_k>T_k$; then $\mathrm{IoU}_{k,c}=\sum_x|M_k\cap L_c|/\sum_x|M_k\cup L_c|$ over the images that have a label of
$c$'s category; unit $k$ detects its best concept if IoU $>0.04$, and a layer's \emph{unique detectors} are the distinct
concepts detected. Intersections for all concepts are accumulated per batch with a GPU \texttt{index\_add} (exact counts),
all sums are kept for 50 random image shards (bootstrap intervals in Fig.~\ref{fig:last}), and the quantile is computed exactly
with bounded memory (VGG conv1\_2 has $3.2\cdot10^9$ activations per unit). A random rotation $Q\in SO(K)$ comes from a QR
decomposition; $Q^\alpha$ from the real Schur form.

\section{Results}

\subsection*{Result 1 --- Interpretability is axis-aligned (paper Sec.~3.2)}
\begin{figure}[t]\centering
\begin{subfigure}{0.5\linewidth}\includegraphics[width=\linewidth]{\fig{fig_rotation.png}}\caption{All rotated models (3 seeds each).}\label{fig:rot}\end{subfigure}\hfill
\begin{subfigure}{0.48\linewidth}\includegraphics[width=\linewidth]{\fig{fig_rotation_places205.png}}\caption{Paper's AlexNet-Places205 (paper Fig.~4).}\end{subfigure}
\caption{Result 1: random rotations $Q^\alpha$ of the last conv layer.}
\end{figure}
For the paper's own model, rotating conv5 by a random $Q\in SO(256)$ reduces unique detectors from 63 to $10.0\pm0.8$
(3 rotations), an \textbf{84\% drop} (paper: 80\%), with the collapse between $\alpha=0.6$ and $0.8$ as in the paper.
VGG16-Places365 conv5\_3: $148\to30.3\pm1.2$ (80\%). The drop is smaller for other models: torchvision AlexNet-ImageNet 58\%,
Caffe AlexNet-ImageNet 68\%, PyTorch AlexNet-Places365 69\% (3 seeds each): ImageNet models drop less than Places models in
both model families. $\alpha=0$ reproduces the unrotated baseline exactly. We also measured the paper's premise: passing
$Q^\top Qf(x)$ through the rest of the network leaves the output unchanged (max.\ logit difference $1.5\cdot10^{-5}$), while
$Qf(x)$ alone gives chance accuracy. Conclusion: interpretability is a property of the basis, not of the information (H2): \ok.

\subsection*{Result 2 --- Concepts by layer (paper Sec.~3.3)}
\begin{figure}[t]\centering
\begin{subfigure}{0.99\linewidth}\centering\includegraphics[width=0.86\linewidth]{\fig{fig_layers_alexnet.png}}\caption{AlexNet, original Caffe models: Places205 (paper's model), ImageNet (BVLC), and a random AlexNet (paper Fig.~5).}\end{subfigure}
\begin{subfigure}{0.99\linewidth}\centering\includegraphics[width=0.99\linewidth]{\fig{fig_layers_vgg.png}}\caption{VGG16, last conv layer of each block.}\end{subfigure}
\caption{Result 2: unique detectors per layer.}\label{fig:layers}
\end{figure}
AlexNet-Places205 rises 7, 20, 34, 37, 63 from conv1 to conv5; colour dominates conv1, textures conv2--conv3, objects conv5
(33 object detectors): \ok. ImageNet and Hybrid AlexNets dip at conv4 (Caffe ImageNet: 35 $\to$ 29, although conv3 and conv4
both have 384 units), so the rise is not monotone for every model. VGG16 stays low up to conv3\_3 (6--13) and jumps at
conv4\_3/conv5\_3, the same pattern for all trained VGGs. Random networks stay low at every layer.

\subsection*{Result 3 --- Architecture and training data (paper Sec.~3.4)}
\begin{figure}[t]\centering
\includegraphics[width=0.97\linewidth]{\fig{fig_last_layer.png}}
\caption{Result 3: last conv layer of all 12 models (paper Fig.~7/8). tv = torchvision, PT = PyTorch release, Caffe = original model.}\label{fig:last}
\end{figure}
VGG16 has 89--148 unique detectors, AlexNet 43--68: \ok. With the original Caffe models Places beats ImageNet clearly
(AlexNet: Places205 63, Places365 68, Hybrid 65 vs ImageNet 43; VGG16: Places365 148 vs ImageNet 89) and the AlexNet order
Places365 $>$ Hybrid $>$ Places205 $>$ ImageNet is the paper's (Fig.~8): \ok. Two caveats: the torchvision ImageNet AlexNet
(53) and the PyTorch Places365 release (55) are almost equal, and the PyTorch release has hundreds of units that are zero on
$>$99.5\% of positions (the Caffe Places365 model has almost none) --- so ``Places $>$ ImageNet'' holds for the original
models, not for every re-implementation. Random networks have 4--6: \ok. Our counts are higher than the paper's (its
Fig.~4/5 axes end at 40, we find 63 for the same AlexNet-Places205); we use 224\,px instead of 227\,px input and cannot
otherwise explain the difference.

\subsection*{Result 4 --- Training conditions and width (paper Sec.~3.5, 3.7)}
\begin{figure}[t]\centering
\begin{subfigure}{0.52\linewidth}\includegraphics[width=\linewidth]{\fig{fig_training_variants.png}}\caption{Variants, 3 seeds each.}\label{fig:train}\end{subfigure}\hfill
\begin{subfigure}{0.46\linewidth}\includegraphics[width=\linewidth]{\fig{fig_training_dynamics.png}}\caption{During training (3 seeds).}\end{subfigure}
\caption{Result 4: small-scale training on Imagenette (paper Fig.~10, 11, 13).}
\end{figure}
The paper trained AlexNet on Places205; we can only train on Imagenette (9.5k images, 60 epochs, 3 seeds per variant), so
these are weak, small-scale tests. \emph{Seeds}: $19.0\pm1.7$ unique detectors at equal accuracy (82.5\%): \ok.
\emph{Iterations}: 4 at initialisation, 4--6 up to iteration 400, 12 at 800, 19 at the end (mean of 3 seeds); object
detectors appear after $\sim$800 iterations: \ok\ (in shape). \emph{No dropout}: $17.3\pm1.5$ (objects 5.0 vs 5.3,
textures 9.7 vs 11.3; Welch $p=0.28$) --- the paper's ``more textures, fewer objects'' is \no. \emph{Batch norm}:
$23.7\pm3.8$, higher than baseline (not significant, $p=0.15$) and accuracy 86.3\%; pretrained VGG16-BN has fewer detectors
than VGG16 in conv3\_3/conv4\_3 but more in conv5\_3 (96 vs 89) --- the paper's ``BN decreases interpretability'' is \no.
\emph{Width}: replacing the FC layers by global average pooling (GAP) alone already changes the result (256 units: $10.3$
unique, 106 detector units, vs 19 / 60 for the baseline), so width must be compared inside the GAP family: 256 $\to$ 768
$\to$ 1024 $\to$ 2048 units give $10.3, 12.3, 13.3, 15.0$ unique detectors and 106, 327, 426, 884 detector units (3 seeds
each). More units give many more detector units but only a few more concepts, and unlike the paper there is no plateau
after 768: \partly.

\subsection*{Result 5 --- Can the labels be trusted? (Molnar's criteria)}
\begin{figure}[t]\centering
\begin{subfigure}{0.48\linewidth}\includegraphics[width=\linewidth]{\fig{fig_cascade.png}}\caption{Cascading randomisation.}\label{fig:casc}\end{subfigure}\hfill
\begin{subfigure}{0.48\linewidth}\includegraphics[width=\linewidth]{\fig{fig_fidelity_alexnet-imagenet.png}}\caption{IoU vs permutation importance, AlexNet-ImageNet conv5.}\label{fig:fid}\end{subfigure}
\caption{Result 5: sanity check and fidelity.}
\end{figure}
\textbf{Sanity check} (cascading randomisation, Fig.~\ref{fig:casc}). Re-initialising only the classifier changes nothing
(labels identical in all three models tested); re-initialising the top conv layer drops the last layer to 5--14 unique
detectors. The method passes, but it describes what a layer represents, not how the network uses it.
\textbf{Fidelity} (Fig.~\ref{fig:fid}). Using permutation feature importance with units as features (Imagenette,
AlexNet-ImageNet conv5), labelled units are \emph{not} more important than unlabelled ones (Mann--Whitney $p=0.80$; VGG16
$p=0.63$) and IoU does not correlate with importance (Spearman $-0.01$); no single unit changes accuracy by more than 0.2\%.
Shuffling all units labelled ``dog'' together (the one test fixed in advance) costs 24.3 points of 1000-way accuracy on the
dog class (13 units; 10{,}000 random groups: mean 3.4, $p=10^{-4}$; VGG16: 29.3 points, $p=0.002$) but under 1 point in
10-way accuracy: the units carry \emph{which breed}, not \emph{dog or not}. For the original Caffe AlexNet-ImageNet the same
test is \emph{not} significant (11.8 points vs random mean 4.6, $p=0.12$): group fidelity of labels depends on the model.

\begin{table}[t]\centering\small
\begin{tabularx}{\linewidth}{@{}lXlp{5.6cm}@{}}\toprule
& Claim & Verdict & Our evidence \\\midrule
1 & Interpretability is axis-aligned; rotation keeps outputs (3.2) & \ok & $-84\%$ for the paper model (paper $-80\%$); 58--84\% across models \\
2 & Colour/texture early, objects late (3.3) & \ok & all trained models; ImageNet AlexNets dip at conv4 \\
3 & VGG $>$ AlexNet; trained $\gg$ random (3.4) & \ok & 89--148 vs 43--68; random 4--6 \\
3 & Places $>$ ImageNet (3.4) & \ok & original models 63--68 vs 43, VGG 148 vs 89; not for torchvision/PyTorch releases \\
3 & Hybrid between Places and ImageNet (3.4) & \ok & AlexNet 65, VGG 147 \\
4 & Seeds converge (3.5) & \ok & $19.0\pm1.7$ \\
4 & Detectors emerge during training (3.5) & \ok & 4 $\to$ 19, objects after $\sim$800 it. \\
4 & No dropout: more texture, fewer objects (3.5) & \no & no difference ($p=0.28$), small scale \\
4 & Batch norm lowers interpretability (3.5) & \no & higher in small models, mixed in VGG16 \\
4 & Wider layer, more detectors (3.7) & \partly & more detector units, few more concepts \\\midrule
5 & Method passes randomisation sanity check & yes & classifier reset: no change; top-conv reset: 5--14 left \\
5 & Labels identify units the network relies on & no / mixed & single units: $p=0.80$; ``dog'' group: $p=10^{-4}$ (tv), $p=0.12$ (Caffe) \\\bottomrule
\end{tabularx}
\caption{Summary of the five results. Of the 10 paper claims tested (Results 1--4), 7 are reproduced, 1 partly and 2 not.}\label{tab:sum}
\end{table}

\begin{table}[t]\centering\small
\begin{tabularx}{\linewidth}{@{}>{\raggedright\arraybackslash}p{4.6cm}X@{}}\toprule
Challenge & Solution \\\midrule
Broden server offline & Internet Archive copy, parallel ranged download; first file was truncated (index reported the compressed size) --- caught by the SHA-1 check, fixed \\
The paper's models exist only as Caffe files & own protobuf/prototxt reader; every model verified by classification before use; a 5-point accuracy gap traced to the resize protocol, not the converter \\
RG-encoded label PNGs, up to 5 overlapping maps, image-level scene/texture labels & decoder that merges maps without double counting; scene/texture as full-image masks; label counts checked against Broden's \texttt{label.csv} \\
Exact top quantile over up to $3.2\cdot10^9$ values per unit; first version used $>$200\,GB RAM on a model with many all-zero units (tie at 0) & level counting + candidate filtering with explicit tie handling, bounded memory; unit test for ties \\
In-place ReLU overwrote tensors captured by hooks & non-in-place ReLU, cloned activations, unit test \\
Upsampling anchored at receptive-field centres & receptive-field start/stride from the layer geometry, \texttt{grid\_sample} \\
GAP head confounds the width experiment & added a 256-unit GAP control \\
Single training runs cannot support conclusions & 3 seeds per variant, Welch tests \\
Our own bugs (e.g.\ wrong weights path in the fidelity run) & found from logs, fixed, re-run; every claim checked against the raw files \\\bottomrule
\end{tabularx}
\caption{Challenges and solutions.}\label{tab:chal}
\end{table}

\FloatBarrier
\section{Challenges and how we solved them}
Table~\ref{tab:chal} lists the main technical problems.

\section{What did not work, bluntly}
Batch norm, dropout and width could only be tested on a 9.5k-image dataset; the paper's effects may need the scale of
Places205, so our ``not reproduced'' means ``not seen at our scale''. Absolute detector counts are at least 1.5$\times$ the
paper's for the same model (its axes end at 40, we find 63), and we cannot fully explain why. ``Places $>$ ImageNet'' and the
rotation effect size depend on which release of a model is used. Self-supervised models (paper Fig.~8--9) and
VGG16-Places205 could not be obtained, so those comparisons are missing.

\section{Conclusion}
The central result holds: interpretable units are a property of the learned basis, and the paper's own model loses 84\% of
its detectors under a random rotation with identical outputs. The comparative results (layer, architecture, training data)
also hold for the original models. The training-regime results do not reproduce at our scale. As an explanation method,
Network Dissection gives reproducible labels for what units respond to, but single labels say little about what the network
relies on --- they should be combined with tests of importance, as Molnar recommends.

\noindent\textbf{Declaration.} \fillin{State the use of AI tools as required by the course policy.}

\begin{thebibliography}{9}\small
\bibitem{bau2017} D. Bau, B. Zhou, A. Khosla, A. Oliva, A. Torralba. Network Dissection: Quantifying Interpretability of Deep Visual Representations. CVPR 2017.
\bibitem{molnar} C. Molnar. \emph{Interpretable Machine Learning}. \url{https://christophm.github.io/interpretable-ml-book}, accessed \fillin{date}.
\end{thebibliography}
\end{document}
